% =====================================================================
%  main_aaai_intro.tex -- rewritten Introduction for AAAI-27 Main Track (MAS)
%  Spine: docs/if_this_paper_disappeared.md (7 points).
%  Target length: ~1.1-1.25 pages two-column.
%  Pour into \section{Introduction} of main_aaai_skeleton.tex.
%
%  Design decisions baked in:
%   - Para 1 hook: human+agent teams inherit loss aversion (why care).
%   - Para 2: the gap -- collapse known but never measured (the unmeasured object).
%   - Para 3: why the supply chain is the ONLY substrate where the real
%     question (where in the network) is askable -- PROMOTED to lead the
%     "why this domain" logic (was 5th reason in main_standard).
%   - Para 4: the move -- boundary as instrument; CPT as reward transform
%     (pre-empts the "reward engineering" kill-shot inline).
%   - Para 5: <<HYPOTHESIS>> verbatim + the contribution ladder preview.
%   - Behavioral number led by human-range lambda~3 (+0.089); lambda=7 demoted.
% =====================================================================

\section{Introduction}
\label{sec:intro}

Cooperative systems increasingly mix human decision-makers with learning agents
trained to act like them. In automated replenishment, fraud review, energy
dispatch, and clinical triage, the deployed agents do not optimise a clean
objective---they inherit human preference structure by design, either because
they are trained on human choices through preference-aligned objectives that are
themselves prospect-theoretic in form~\citep{ethayarajh2024kto}, or because they
act alongside humans in the loop~\citep{lalmohammed2025}. The most robustly
documented feature of that preference structure is \emph{loss aversion}:
people---and now the agents shaped after them---weigh a loss more heavily than an
equivalent gain~\citep{kahneman1979,tversky1992}. Individual-level behavioral
economics tells us what this does to one decision-maker. It does not tell us what
it does to the \emph{coordination} of many decision-makers learning
together~\citep{leibo2017multi,lowe2017multi}---which is the question a deployer
actually faces.

That coordination has a well-documented way of failing. Cooperative multi-agent
reinforcement learning (MARL) collapses into a degenerate, zero-action
\emph{lazy-agent} policy when agents abandon cooperation~\citep{liu2023lazy,peng2023curo}.
Whether a system collapses depends on its \emph{reward centralization}
$w\in[0,1]$---the fraction of each agent's reward drawn from the team's global
outcome rather than its own local one: push $w$ too low and each agent's own
performance drowns out its contribution to the team, so cooperation carries no
learnable signal and the safe, lazy policy wins. Yet the field describes this
dependence only qualitatively---``centralization helps''~\citep{dmarl2025,liu2023lazy}.
\emph{Where}, in $w$, coordination gives way to collapse---the critical
centralization $\wcrit$ at which the collapse probability crosses one-half---has
never been measured. It is a slogan, not a quantity.

We ground the measurement in automated replenishment for multi-echelon supply
chains, and the choice is not incidental: it is the kind of domain where the
sharpest version of our question can even be \emph{asked}. A supply chain has a
\emph{directed flow}---inventory moves downstream along a chain, so ``upstream''
and ``downstream'' are well defined---which lets us ask not merely whether a
behavioral bias matters but \emph{where in the network} a biased participant
matters most. Its failure modes are economically legible (hoarding,
stockouts, the bullwhip cascade), so collapse is a recognisable operational
disaster rather than an abstract training pathology; and real replenishment is
done by planners who demonstrably weigh losses more heavily than gains, so the
agents replacing or assisting them inherit that bias by construction. The
substrate thus supplies the directed structure, the legible failure, and the
behavioral mechanism in one place---which is why our most fine-grained results
live here, and why we are explicit about which of them we claim beyond it.

We treat collapse the way physics treats a phase boundary: a pre-registered,
frozen binary collapse definition, logistic estimation of
$P(\mathrm{collapse}\mid w)$, and bootstrap confidence intervals on $\wcrit$. We
use ``phase boundary'' in the empirical sense---a sharp, metastable threshold
whose asymptotic ($T\!\to\!\infty$) value is the boundary proper---not as a claim
of a thermodynamic-limit singularity. The boundary is then an \emph{instrument}.
The decisive design choice that makes what it measures trustworthy is that we
apply cumulative prospect theory (CPT) as a \emph{reward transform after reward
blending}, never inside the learning update~\eqref{eq:reward}: the optimiser is
byte-for-byte identical across all conditions, so a measured shift in $\wcrit$ is
a behavioral effect, not an artifact of tuning or of a changed
algorithm---separating us cleanly from prior CPT-in-MARL work that folds the
value and probability-weighting functions into the actor--critic
update~\citep{lalmohammed2025}. Against this null, a three-way design
(no PT / curvature only / curvature plus loss aversion) attributes any shift in
$\wcrit$ to a specific behavioral component.

What we find is this.
% <<HYPOTHESIS>> -- verbatim, occurrence 2 of 3 (also in abstract and conclusion)
\emph{Human-like loss aversion forces cooperative learning agents to require more
reward centralization to avoid collapse---a shift we predict in sign a priori and
confirm across three environments and two algorithms---and a single downstream
loss-averse agent can set that requirement for the whole system.} The shift is
already significant at the human-range $\lambda\!\approx\!3$
($+0.089$ $[0.041,0.141]$ in required centralization), rising to a high-stakes
upper bound of $+0.166$ $[0.135,0.202]$ at $\lambda{=}7$; its \emph{sign} is fixed
before any data by a sign invariant of the policy-gradient escape barrier
(\Cref{thm:sign}); and its \emph{direction} replicates across a supply chain
(under PPO and centralized-critic MAPPO), MPE Cooperative Navigation, and
Level-Based Foraging. Our contributions form a ladder: (i) the first
quantitative, confidence-intervalled phase boundary of coordination collapse in
deep cooperative MARL; (ii) the discovery that this boundary is \emph{metastable},
sliding with the training horizon rather than being a fixed constant; (iii) a
behavioral decomposition showing the two CPT components move it in \emph{opposite}
directions, with the loss-aversion direction generalizing and a \emph{minority by
network position} setting the boundary; (iv) an a priori sign guarantee for that
direction; and (v) a practical reading of the boundary as a measurable safety
margin---how much more centralization a cooperative system needs once its agents
are loss-averse, and which participant sets the requirement. Two of these results
(the affine $\lambda$-law and the downstream-minority cascade) we establish on the
supply-chain substrate and scope honestly as such, framing the cascade as a
falsifiable, structurally specified prediction for any directed-flow domain.
